\scanpage{57}
\mspage{\unsure{44|41}}
% archive sheet number: 56

puis que le carré ayant deux côtés $u_{K,J},u'_{K,J}$ est cocartésien. (NB Comparer le diagramme similaire p.~32 -- il y aurait un lemme commun à dégager, distinguant \unsure{deux|des} cas \ldots)

\deleted{Pour prouver que} ($\mathrm{Cof}_{\unsure{2|1}}$) Pour prouver que \deleted{$\unsure{f}$} $C_{I^{\unsure{\circ|0}}}\cap W_{I^\circ}$ est stable par cochangement de base, je ne vois pas de façon de le faire, si je ne sais déjà que

\[
 C_{I^\circ}\cap W_{I^\circ}=(TC)_{I^\circ},\tag*{(\unsure{\alpha})}
\]

et alors on applique \unsure{le} \deleted{\unsure{vérificat}} résultat précédent (où $W$ n'entrait pas) à $TC$ \unsure{ald} $C$. Mais pour avoir \unsure{ceci|ces}, il faut que je suppose qu'on est dans l'un des deux cas suivants (de façon à pouvoir appliquer le th.~2.1. et son corollaire 2.5.)

a) $I$ fini, ou plus généralement $I$ \unsure{artinien} \uline{et} $\mathrm{Top}(I)$ est noethérien

b) $I$ artinien, et de plus $C$ et $TC$ stables par composition transfinie.

($\mathrm{Cof}_3$) L'axiome de factorisation me paraît \unsure{à|a} priori le plus délicat$^{\times}$. Mais cela \deleted{\unsure{tout}} fonctionne sans trop de mal. Soit donné
\[
 f:X_{\unsure{\cdot|\bullet}}\unsure{\mathbin{-}|\longrightarrow}Y_{\cdot}.
\]

% slanted
\margin{$^{\times}$L'argument qui suit s'applique \unsure{à d'autres} \unsure{situations} de factorisation, par exemple pour une str. \unsure{proprement} de Quillen, donnée $(\unsure{\Phi,\Psi|C,F})$ dans $M$.}
